\documentclass{ieeeaccess}
\usepackage{cite}
\usepackage{amsmath,amssymb,amsfonts}
\usepackage{graphicx}
\usepackage{textcomp}
\usepackage{booktabs}
\usepackage{multirow}
\usepackage{url}
% Lightweight pseudocode environment using only base packages.
\newcounter{algocnt}
\newenvironment{savpseudocode}[1]
  {\refstepcounter{algocnt}\par\noindent
   \begin{minipage}{\columnwidth}\hrule\smallskip
   \textbf{Algorithm \arabic{algocnt}.\ #1}\smallskip\hrule\smallskip
   \begin{tabbing}
   xx\=xx\=xx\=xx\=xx\=xx\=\kill}
  {\end{tabbing}\hrule\end{minipage}\smallskip}

\usepackage{bm}
\makeatletter
\AtBeginDocument{\DeclareMathVersion{bold}
\SetSymbolFont{operators}{bold}{T1}{times}{b}{n}
\SetSymbolFont{NewLetters}{bold}{T1}{times}{b}{it}
\SetMathAlphabet{\mathrm}{bold}{T1}{times}{b}{n}
\SetMathAlphabet{\mathit}{bold}{T1}{times}{b}{it}
\SetMathAlphabet{\mathbf}{bold}{T1}{times}{b}{n}
\SetMathAlphabet{\mathtt}{bold}{OT1}{pcr}{b}{n}
\SetSymbolFont{symbols}{bold}{OMS}{cmsy}{b}{n}
\renewcommand\boldmath{\@nomath\boldmath\mathversion{bold}}}
\makeatother

\def\BibTeX{{\rm B\kern-.05em{\sc i\kern-.025em b}\kern-.08em
    T\kern-.1667em\lower.7ex\hbox{E}\kern-.125emX}}

% Placeholder for numbers that will be filled in once the iDRAC remote
% experiment grid completes.  Grep for \TBD to find every empirical hole.
\newcommand{\TBD}[1]{\textbf{[TBD: #1]}}

\begin{document}
\history{Manuscript prepared April 2026; submitted to IEEE Access.}
\doi{10.1109/ACCESS.2026.XXXXXXX}

\title{Safety-Aware Vehicle Routing for Quick-Commerce Last-Mile
Delivery in Indian Cities: An $\varepsilon$-Constraint Multi-Objective
Formulation with Calibrated Crash Risk}

\author{\uppercase{Subhajit Sidhanta}\authorrefmark{1},
\uppercase{Pranav Bhadane}\authorrefmark{1}}

\address[1]{Department of Industrial and Systems Engineering,
Indian Institute of Technology Kharagpur, Kharagpur, West Bengal 721302, India
(e-mail: subhajit@iem.iitkgp.ac.in;
pranavbhadane@kgpian.iitkgp.ac.in)}

\tfootnote{This work was carried out at the Department of Industrial
and Systems Engineering, Indian Institute of Technology Kharagpur.
All software, configuration, and the BASM calibration note are
released under the MIT licence at the project repository.}

\markboth
{Sidhanta \& Bhadane: Safety-Aware VRPTW for Quick-Commerce Delivery in Indian Cities}
{Sidhanta \& Bhadane: Safety-Aware VRPTW for Quick-Commerce Delivery in Indian Cities}

\corresp{Corresponding author: Subhajit Sidhanta
(phone: +91-3222-284750; e-mail: subhajit@iem.iitkgp.ac.in).}

\begin{abstract}
Quick-commerce platforms in India promise ten-minute delivery from
neighbourhood ``dark stores'' to customers, putting persistent time
pressure on a fleet of two-wheeler riders that must operate within an
already-hazardous urban traffic environment. Conventional vehicle
routing formulations minimise travel time and treat safety as a
post-hoc concern. We propose a Safety-Aware Vehicle Routing Problem
with Soft Time Windows (SA-VRPTW) that explicitly bounds per-route
crash exposure and fleet-wide use of vulnerable residential roads via
$\varepsilon$-constraints, while minimising a convex soft-time-window
penalty as the primary objective. The crash risk per road segment is
calibrated through a documented public-source chain
(MoRTH 2022 + Mohan~et~al.~2017 + an OpenStreetMap structural proxy)
that does not require access to the protected iRAD microdata. We
implement three solvers under a unified evaluation contract: a
piecewise-linear MILP, a genetic algorithm with Bellman--Ford optimal
splitting, and an Adaptive Large Neighbourhood Search with
Ropke--Pisinger adaptive operator weights. We benchmark on five Indian
cities (Bengaluru, Delhi, Gurugram, Mumbai, Pune) with real Blinkit
dark-store coordinates extracted from a public KML snapshot, sweeping
instance sizes $N \in \{20, 50, 100, 200\}$ over ten seeds. The
empirical study traces Pareto frontiers between $F_1$ delivery
performance and per-route crash survival, quantifies the cost of the
residential-road constraint, and verifies behavioural sensitivity
through a Beta-prior compliance Monte Carlo. The contribution is a
fully reproducible, calibration-honest framework that lets a
quick-commerce operator dial the safety/efficiency trade-off
explicitly rather than implicitly.
\end{abstract}

\begin{keywords}
Adaptive large neighbourhood search,
crash risk modelling,
$\varepsilon$-constraint multi-objective optimisation,
genetic algorithms,
last-mile delivery,
mixed-integer linear programming,
quick-commerce,
soft time windows,
urban routing,
vehicle routing problem.
\end{keywords}

\titlepgskip=-21pt

\maketitle

\section{Introduction}
\label{sec:introduction}
\PARstart{Q}{uick-commerce} (q-commerce) platforms such as Blinkit,
Zepto, and Instamart have rewired Indian urban delivery by promising
arrival times measured in minutes rather than hours. The operating
pattern is a network of small ``dark stores'' (also called depots or
forward-deployed micro-warehouses), each serving a $\sim$2--3~km radius
through a fleet of two-wheeler riders. Order batches are dispatched
every few minutes, with each rider typically carrying one to two
items per route to keep the delivery promise realistic. The combination
of strict service-level agreements (SLAs), high dispatch frequency,
narrow residential-street penetration, and a vulnerable rider
population (helmeted but exposed two-wheelers in mixed traffic) creates
a routing problem in which safety is not an optional refinement.

The classical Vehicle Routing Problem with Time Windows (VRPTW)
formulation captures none of this. It minimises a weighted sum of
travel time, distance, or vehicle count subject to capacity and time
windows. Crash risk, when modelled, is typically embedded as an
additional weighted term in the objective using a constant
``risk-per-edge'' coefficient that has no calibration story. Several
works in safety-aware routing penalise hazardous arcs implicitly via
inflated travel times, which conflates two distinct phenomena
(slow vs.\ dangerous) and gives the operator no explicit knob.

The contribution of this paper is a Safety-Aware VRPTW (SA-VRPTW)
formulation that:
\begin{enumerate}
\item Treats crash exposure as a hard $\varepsilon$-constraint per route
(survival probability $\geq e^{-\bar{R}}$) rather than as an objective
component, so the safety/efficiency trade-off is traced
\emph{explicitly} as a Pareto frontier;
\item Adds a fleet-wide budget on the use of vulnerable residential
streets, so a planner can cap network-level penetration of narrow
neighbourhood roads even when individual routes are safe;
\item Calibrates per-edge crash probability through a fully documented
public-source chain (MoRTH road-category fatality totals
\cite{morth2022} + Mohan~et~al.~2017
\cite{mohan2017} ordering constraints
+ an OpenStreetMap structural proxy), avoiding the need for the
protected Integrated Road Accident Database (iRAD) microdata;
\item Replaces the linear lateness term commonly used in soft-time-window
formulations with a convex exponential SLA penalty
$\exp(\beta \tau) - 1$, which we approximate via piecewise-linear
outer-envelope in the MILP and evaluate exactly in the heuristics;
\item Is implemented behind a unified solver/evaluation contract so
that an exact MILP, a giant-tour-with-Bellman--Ford-split genetic
algorithm, and an adaptive large neighbourhood search all evaluate
solutions through identical $F_1$ and feasibility code, eliminating
the ``apples-to-pears'' bias that plagues many VRP comparisons;
\item Uses real Blinkit dark-store coordinates extracted from a
publicly committed KML snapshot, not synthetic uniform-disk depots, so
the geometric difficulty of the instances reflects the actual
topology operators face.
\end{enumerate}

The remainder of the paper is organised as follows. Section~\ref{sec:related}
surveys related work in safety-aware routing and quick-commerce
optimisation. Section~\ref{sec:formulation} states the SA-VRPTW
formulation in full. Section~\ref{sec:calibration} describes the BASM
calibration chain, the BPR congestion model, and the dark-store data
pipeline. Section~\ref{sec:methods} details the three solvers.
Section~\ref{sec:experiments} reports the empirical study.
Section~\ref{sec:discussion} discusses operator-facing implications,
limitations, and threats to validity. Section~\ref{sec:conclusion}
concludes.

\section{Related Work}
\label{sec:related}

\subsection{VRPTW with Soft Time Windows}
The classical hard-time-window VRPTW \cite{solomon1987} requires every
customer to be served strictly inside a window $[e_i, l_i]$. The soft
variant \cite{koskosidis1992,fu2008} permits late arrival at a
penalty, typically linear in lateness. Convex penalties have been
proposed for SLA-driven settings \cite{taillard1997} but their
linearisation in MILP models is non-trivial and is usually omitted.
We adopt $\exp(\beta \tau) - 1$ to reflect the contractual cost
structure of q-commerce, where lateness becomes operationally
expensive faster than linearly (rider re-routing, customer churn,
SLA refunds).

\subsection{Safety-Aware Vehicle Routing}
Hoseinzadeh~et~al.\ \cite{hoseinzadeh2020} model cumulative crash risk
along a tour but use weighted-sum aggregation, which makes the
safety/efficiency trade-off implicit and difficult to interpret.
$\varepsilon$-constraint methods \cite{salmon2023,ehrgott2005} dominate
weighted-sum approaches when the planner needs to read off a Pareto
frontier rather than a single weighted optimum, and we adopt them
here.

\subsection{Quick-Commerce and Last-Mile Routing}
Sinchaisri~et~al.\ \cite{sinchaisri2023} document the operational
intensity of q-commerce dispatch in the United States. Field studies
of Indian gig-delivery rider behaviour
\cite{das2020,zografos2004} highlight a tension between SLA pressure
and safe driving compliance. Our behavioural Monte Carlo
(Section~\ref{subsec:behavioural}) is calibrated against these
findings, with a sensitivity sweep documenting how the conclusions
move as the assumed compliance prior shifts.

\subsection{Crash-Risk Calibration in Routing}
Most safety-aware routing papers treat the per-edge crash probability
as exogenous, with calibration relegated to a single citation or
synthetic distribution. We take a different stance: every constant
in our calibration file is traced to a public source, and we publish
the derivation as a separate calibration note (Section~\ref{sec:calibration}).

\section{Problem Formulation}
\label{sec:formulation}

\subsection{Notation}
Let $G = (V, E)$ be the directed street graph for a city, extracted
via OSMnx \cite{boeing2017osmnx} with the drive network filter. Let $D \subset V$ be the set
of depot nodes (k-means cluster centroids of the city's Blinkit dark
stores, snapped to the nearest OSM nodes), and $C \subset V$ the set
of customer nodes (sampled with residential-density weighting,
constrained to a delivery radius around depots). Let
$K = \bigcup_{d \in D} K_d$ be the rider fleet, partitioned by home
depot.

For each street-graph edge $(i, j) \in E$:
\begin{itemize}
\item $\ell_{ij}$ is the length in metres,
\item $v_{ij}^{\text{free}}$ is the free-flow speed (km/h, OSMnx imputed),
\item $t_{ij}^{\text{free}} = 60 \ell_{ij} / (1000 v_{ij}^{\text{free}})$ is the free-flow travel time in minutes,
\item $\rho_{ij} = \alpha (V_{ij}/C_{ij})^{\beta}$ is the Bureau of Public Roads (BPR) \cite{bpr1964} congestion factor with $\alpha = 0.15$, $\beta = 4$,
\item $t_{ij} = t_{ij}^{\text{free}} (1 + \rho_{ij})$ is the congestion-inflated travel time used everywhere downstream,
\item $r_{ij} \in [0, 0.99]$ is the per-traversal crash probability (BASM, Section~\ref{subsec:basm}),
\item $h_{ij} \in \{0, 1\}$ indicates a residential street (OSM \texttt{highway} $\in$ \{\texttt{residential}, \texttt{living\_street}\}).
\end{itemize}

For each pair $(i, j)$ with $i, j \in C \cup D$, the super-graph arc
$(i, j)$ is the shortest-time path on $G$, with cumulative quantities
\begin{align}
T_{ij} &= \sum_{e \in P_{ij}} t_e, \label{eq:Tij} \\
R_{ij} &= \sum_{e \in P_{ij}} -\ln(1 - r_e), \label{eq:Rij} \\
H_{ij} &= \sum_{e \in P_{ij}} h_e. \label{eq:Hij}
\end{align}
$R_{ij}$ is additive in log-survival because per-traversal crashes are
treated as independent events along the path.

\subsection{Customer Parameters}
Each customer $i \in C$ has a demand $q_i \in \{1, 2\}$, an order
placement time $e_i \in [0, 5]$ minutes, a promised arrival
$\mathrm{ETA}_i = e_i + 10$ (the q-commerce ten-minute promise), and a
service time $s_i \in \{2, 4\}$ minutes (4~min for high-rise
buildings, otherwise 2). Time windows are soft.

\subsection{Decision Variables}
For rider $k \in K$ and arc $(i, j)$ in the super-graph:
\begin{itemize}
\item $x_{ij}^k \in \{0, 1\}$: rider $k$ traverses arc $(i, j)$,
\item $a_i^k \geq 0$: arrival time of rider $k$ at node $i$ (defined only when $i$ is visited),
\item $w_i^k \geq 0$: idle wait at $i$ (early arrival),
\item $\tau_i^k \geq 0$: lateness at $i$,
\item $u_i^k \in \mathbb{Z}_{\geq 1}$: MTZ sequence variable for sub-tour elimination.
\end{itemize}

\subsection{Objective}
The primary objective $F_1$ minimises ETA-deviation plus convex SLA
lateness penalty across the fleet:
\begin{equation}
F_1 = \sum_{k \in K} \sum_{i \in C}
\left[ w_{\text{early}} \cdot w_i^k + \big(\exp(\beta_{\text{stw}} \tau_i^k) - 1\big) \right].
\label{eq:F1}
\end{equation}
Travel time enters implicitly through the arrival-time linking
constraints; congestion is already baked into $t_{ij}$; crash risk is
handled as an $\varepsilon$-constraint, not as an objective term, so
the operator can read its cost off the Pareto frontier.

\subsection{$\varepsilon$-Constraints}
Per-route crash-survival budget:
\begin{equation}
\sum_{(i,j)} R_{ij} x_{ij}^k \leq \bar{R}, \quad \forall k \in K.
\label{eq:Rbar}
\end{equation}
Fleet-wide residential-edge budget:
\begin{equation}
\sum_{k \in K} \sum_{(i,j)} H_{ij} x_{ij}^k \leq \bar{H}.
\label{eq:Hbar}
\end{equation}
Per-route residential cap:
\begin{equation}
\sum_{(i,j)} H_{ij} x_{ij}^k \leq \bar{H}_k, \quad \forall k \in K.
\label{eq:Hk}
\end{equation}
The pair $(\bar{R}, \bar{H})$ is swept on a $7 \times 6$ grid
(Section~\ref{sec:experiments}) to trace the Pareto frontier; non-dominated
points are retained.

\subsection{Core VRPTW Constraints}
Each customer visited exactly once:
\begin{equation}
\sum_{k \in K} \sum_{j \in C \cup D, j \neq i} x_{ji}^k = 1,
\quad \forall i \in C. \label{eq:visit}
\end{equation}
Flow conservation at every customer:
\begin{equation}
\sum_{i} x_{ij}^k = \sum_{i} x_{ji}^k, \quad \forall j \in C, \forall k. \label{eq:flow}
\end{equation}
Capacity:
\begin{equation}
\sum_{i \in C} q_i \sum_{j} x_{ij}^k \leq Q, \quad \forall k. \label{eq:cap}
\end{equation}
Multi-depot coupling: rider $k$ home-depoted at $d$ leaves $d$ at most
once, returns iff it left, and cannot touch any other depot (omitted
for brevity; full equations in the open-source repository).

Sub-tour elimination (MTZ):
\begin{equation}
u_i^k - u_j^k + |C| x_{ij}^k \leq |C| - 1, \quad \forall i, j \in C, \forall k. \label{eq:mtz}
\end{equation}
Arrival-time linking:
\begin{equation}
a_j^k \geq a_i^k + s_i + T_{ij} - M (1 - x_{ij}^k), \quad \forall i, j, k. \label{eq:link}
\end{equation}
Maximum route duration:
\begin{equation}
\sum_{(i,j)} T_{ij} x_{ij}^k + \sum_{i} s_i \sum_j x_{ji}^k \leq T_{\max}, \quad \forall k. \label{eq:Tmax}
\end{equation}
Soft time windows:
\begin{align}
w_i^k &\geq e_i - a_i^k, \quad w_i^k \geq 0, \label{eq:early} \\
\tau_i^k &\geq a_i^k - \mathrm{ETA}_i, \quad \tau_i^k \geq 0. \label{eq:late}
\end{align}

\subsection{MILP Linearisation of the SLA Penalty}
The MILP cannot evaluate $\exp(\beta_{\text{stw}} \tau)$ directly. We
approximate via a piecewise-linear outer envelope at breakpoints
$\tau \in \{0, 2, 5, 10, 15, 20, 30\}$ minutes, encoded with SOS2 sets.
The heuristics evaluate the exponential exactly. Both share the
same feasibility checker and report $F_1$ through the same
\texttt{savrptw.eval.objective} module.

Constants: $Q = 2$, $T_{\max} = 35$~min,
$\bar{H}_k = 100$ residential edges per route,
$\beta_{\text{stw}} = 0.12$, $w_{\text{early}} = 1.0$, $M = 10^4$.
Per-depot rider counts are sized as
$K_d = \lceil 1.2 N_d / Q \rceil$, providing a 20\% over-provision
buffer above the no-slack lower bound.

\section{Calibration and Data Pipeline}
\label{sec:calibration}

\subsection{Bengaluru Accident Surrogate Model (BASM)}
\label{subsec:basm}

The BASM converts publicly documented aggregate Indian road-safety
rates into a transparent edge-level prior, which is then redistributed
within each OSM road class using graph-local structure. The functional
form is
\begin{equation}
r_{ij} = \mathrm{clip}\!\left[
\lambda_{\text{class}} \cdot \frac{\ell_{ij}}{1000} \cdot
\sigma_{\text{class}} \cdot \pi_{ij},
0, 0.99 \right],
\label{eq:BASM}
\end{equation}
where:
\begin{itemize}
\item $\lambda_{\text{class}}$ is a per-km, per-traversal base
probability for the OSM road class, derived from MoRTH 2022 fatal+grievous
events per category-km divided by the project's BPR-derived annual
traversal exposure;
\item $\sigma_{\text{class}}$ is a rank-preserving severity multiplier
(arterial $>$ collector $>$ local), computed via the bounded transform
$\sigma = \sqrt{\text{intensity} / g}$ where intensity is the
fatality-share-to-length-share ratio per category and $g$ is the
geometric mean of intensities;
\item $\pi_{ij} \in [0.5, 2.0]$ is a bounded local proxy combining
approximate edge betweenness, traffic-signal density, and crossing
density at the edge endpoints.
\end{itemize}
The detailed derivation, including the page references for every
MoRTH number used, is documented separately in the project's
\texttt{BASM\_CALIBRATION.md} note. We do not claim access to the
protected iRAD microdata.

\subsection{BPR Congestion Model}
\begin{equation}
\rho_{ij} = 0.15 \left(\frac{V_{ij}}{C_{ij}}\right)^4,
\label{eq:BPR}
\end{equation}
with capacity $C_{ij}$ from HCM~2010 \cite{hcm2010} Exhibit~10-3
indexed by \texttt{(highway\_class, lane\_count)} and flow $V_{ij}$
derived from a
class-level AADT prior multiplied by an hour-of-day factor (peak
windows 08--11 and 18--21). All values are in
\texttt{conf/congestion/bpr.yaml}.

\subsection{Dark Stores}
We use the publicly available KML snapshot at
\url{https://darkstoremap.in/dark_store.kml}, committed to the
repository as
\texttt{data/raw/darkstoremap\_in\_2026-04-09.kml}. The committed
snapshot contains 461 Blinkit placemarks; per-city
in-bounding-box counts are summarised in
Table~\ref{tab:darkstore-counts}. Hyderabad is excluded from the
primary single-brand experiments because the snapshot has zero in-bbox
Blinkit placemarks there; refetching on 2026-04-25 produced a
byte-identical file. For each retained city, we apply k-means
clustering with $k \in \{3, 5, 8\}$ depots to obtain depot centroids.

\begin{table}
\caption{\textbf{Blinkit dark-store coverage per city in the
2026-04-09 KML snapshot.}}
\label{tab:darkstore-counts}
\centering
\begin{tabular}{lc}
\toprule
City & In-bbox count \\
\midrule
Bengaluru & 113 \\
Delhi & 111 \\
Mumbai & 96 \\
Pune & 39 \\
Gurugram & 38 \\
Hyderabad & 0 (excluded) \\
\bottomrule
\end{tabular}
\end{table}

\subsection{Cross-Validation Against OpenStreetMap}
We attempted to cross-validate the KML coordinates against
OpenStreetMap by querying the Overpass API for features tagged
\texttt{brand=Blinkit} or \texttt{name~/Blinkit/i} within radii of
250~m and 1~km of each sampled coordinate. At a 1~km radius across
24~stores, only 2 matched (8\%, 95\% Wilson CI [0.02, 0.26]), driven
by a single hit in Delhi and one in Gurugram. We interpret this as
evidence about OSM's sparse coverage of Indian dark stores rather
than about KML accuracy, and document this conclusion in the
repository's \texttt{DARKSTORE\_VALIDATION.md} note.

\section{Solution Methods}
\label{sec:methods}

We implement three solvers behind a uniform contract: each consumes a
typed \texttt{Instance} dataclass and returns a typed \texttt{Solution}.
All three evaluate solutions through the same
\texttt{savrptw.eval.objective.breakdown()} and
\texttt{savrptw.eval.feasibility.validate()} functions, ensuring that
reported $F_1$ values are strictly comparable.

\subsection{MILP}
A PuLP/CBC model encodes the formulation in
Section~\ref{sec:formulation} directly. Two arrival-variable groups
are kept separate to avoid a depot-arrival overload: $a_i^k$ for
customer arrivals and $a_{\text{return}}^k$ for the home-depot return,
which prevents the dispatch ($a_d^k = 0$) and return ($a_d^k > 0$)
constraints from contradicting each other. The convex SLA penalty is
linearised by SOS2-encoded outer-envelope segments at the breakpoints
listed above. CBC is given a 120~s wall-clock budget per row in the
default benchmark grid; longer budgets are needed only for $N \geq 50$
where the search is genuinely difficult.

\subsection{Genetic Algorithm}
The GA uses the giant-tour-with-Bellman--Ford-split decoder of
Prins~\cite{prins2004}: each chromosome is a permutation over $C$;
the decoder finds the optimal partition into routes by solving a
shortest-path problem on the auxiliary acyclic graph whose arcs
correspond to feasible single-route cuts. Operators are OX1 crossover
plus a randomised mutation choice over swap, insertion, and segment
reversal (the last subsumes the 2-opt move). Tournament selection
($k = 5$), elitism (top 4 retained), default population 100,
200~generations.

\subsection{Adaptive Large Neighbourhood Search}
The ALNS implementation follows
Ropke and Pisinger~\cite{ropke2006}. Destroy operators include random
removal, worst removal (largest insertion-cost contribution), Shaw
removal (relatedness over travel-time proximity), and a custom
risk-cluster removal that targets customers in the highest-$R_{ij}$
neighbourhoods. Repair operators include cheapest-feasible insertion,
regret-2, and regret-3. Operator weights are updated in segments of
100~iterations using the Ropke--Pisinger reaction factor; rewards
are $\sigma = (33, 9, 1, 0)$ for new-best, improving, accepted-worse,
and rejected moves. Acceptance is simulated annealing with a
geometric cooling rate of 0.99975 starting from a temperature equal
to 5\% of the initial $F_1$. Algorithm~\ref{alg:alns} states the full
loop.

\begin{savpseudocode}{ALNS for SA-VRPTW}\label{alg:alns}
\textbf{Input:}\> Instance $\mathcal{I}$; iterations $T$; segment length $L$; \\
\>\> reaction factor $r$; reward weights $\sigma$; \\
\>\> cooling rate $\eta$; init.\ temperature ratio $\theta$ \\
\textbf{Output:}\> Best feasible solution $s^\star$ \\[2pt]
\textbf{1.}\> $s_0 \gets \mathrm{GreedyConstruct}(\mathcal{I})$ \\
\textbf{2.}\> $s \gets s_0$;\ $s^\star \gets s_0$;\ $F^\star \gets F_1(s_0)$ \\
\textbf{3.}\> $T_{\mathrm{SA}} \gets \theta \cdot F^\star$ \\
\textbf{4.}\> $w^d_o, w^r_o \gets 1$ for every operator $o$ \\
\textbf{5.}\> \textbf{for} $t = 1$ \textbf{to} $T$ \textbf{do} \\
\textbf{6.}\>\> $q \gets \lceil n \cdot \mathcal{U}(0.15, 0.45) \rceil$ \\
\textbf{7.}\>\> $D \gets$ destroy operator sampled $\propto w^d$ \\
\textbf{8.}\>\> $R \gets$ repair operator sampled $\propto w^r$ \\
\textbf{9.}\>\> $s' \gets R(D(s, q))$;\ $\Delta \gets F_1(s') - F_1(s)$ \\
\textbf{10.}\>\> \textbf{if} $F_1(s') < F^\star$ \textbf{then} \\
\>\>\> $s^\star \gets s'$;\ $F^\star \gets F_1(s')$;\ $\rho \gets \sigma_1$ \\
\textbf{11.}\>\> \textbf{else if} $\Delta < 0$ \textbf{then} $\rho \gets \sigma_2$ \\
\textbf{12.}\>\> \textbf{else if} $\mathcal{U}(0,1) < \exp(-\Delta/T_{\mathrm{SA}})$ \textbf{then} $\rho \gets \sigma_3$ \\
\textbf{13.}\>\> \textbf{else} $\rho \gets \sigma_4$ \\
\textbf{14.}\>\> \textbf{if} $\rho > \sigma_4$ \textbf{then} $s \gets s'$ \\
\textbf{15.}\>\> $\pi^d_D \mathrel{+}= \rho$;\ $\pi^r_R \mathrel{+}= \rho$ \\
\textbf{16.}\>\> \textbf{if} $t \bmod L = 0$ \textbf{then} \\
\>\>\> $w^*_o \gets (1{-}r)\, w^*_o + r\, \pi^*_o / n^*_o$ \\
\>\>\> reset $\pi^*$ and counts \\
\textbf{17.}\>\> $T_{\mathrm{SA}} \gets \eta \cdot T_{\mathrm{SA}}$ \\
\textbf{18.}\> \textbf{return} $s^\star$
\end{savpseudocode}

\subsection{Crash-Survival Monte Carlo}
For evaluation only (not optimisation), each route is sampled
$10\,000$ times: each edge traversal is a Bernoulli trial with success
probability $1 - r_e$, and the route survives iff every edge
traversal succeeds. We report mean route survival, fleet survival, and
Wilson 95\% confidence intervals \cite{wilson1927}. The analytic per-route survival
$\exp(-R_{\text{route}})$ is reported alongside as a sanity check.

\subsection{Behavioural Compliance Sensitivity}
\label{subsec:behavioural}
Rider compliance is modelled as a
$\mathrm{Beta}(\alpha, \beta)$ prior. The central scenario is
$(\alpha, \beta) = (5, 2)$. We sweep
$\alpha \in \{3, 5, 8\}$ and $\beta \in \{1, 2, 3\}$, yielding nine
operating points. For each, we Monte-Carlo a $50$-rider population over
$5\,000$ trips and report the IQR band on per-rider survival, with
edge-level traversal rates inflated by a factor
$1 + k_c (1 - c)$ where $c$ is the sampled compliance level
($k_c = 0.4$ matched to Indian gig-rider field data
\cite{das2020}).

\section{Experimental Study}
\label{sec:experiments}

\subsection{Setup}
We benchmark on five Indian cities (Bengaluru, Delhi, Gurugram,
Mumbai, Pune) with depot counts $\{3, 5, 8\}$ scaled to $N$ as in
Table~\ref{tab:depot-N}. Customer counts are
$N \in \{20, 50, 100, 200\}$. Each (city, $N$, solver) cell is
replicated over $10$ random seeds. The MILP is capped at $N \leq 50$
because CBC does not scale gracefully past that point inside the
120~s budget. Solvers are run on \TBD{remote iDRAC compute node
specification (CPU model, RAM, OS)}.

\begin{table}
\caption{\textbf{Per-instance depot counts as a function of $N$.}}
\label{tab:depot-N}
\centering
\begin{tabular}{cc}
\toprule
$N$ & depots \\
\midrule
20 & 3 \\
50 & 5 \\
100 & 5 \\
200 & 8 \\
\bottomrule
\end{tabular}
\end{table}

\subsection{Reproducibility}
The full pipeline (OSM graph $\rightarrow$ BPR congestion $\rightarrow$
calibrated risk $\rightarrow$ Blinkit dark-store clustering $\rightarrow$
super-graph construction $\rightarrow$ solver $\rightarrow$ evaluator)
is composable through Hydra. A complete experiment row is reproducible
from
\begin{verbatim}
PYTHONPATH=src python scripts/run_experiment.py \
  city=bengaluru solver=ga instance.N=50 \
  instance.seed=42 instance.R_bar=0.40
\end{verbatim}
and the full grid through
\texttt{scripts/run\_grid.py}, which checkpoints into a manifest so
remote runs can resume from a crash without recomputing. Every
parameter has a single source-of-truth file in \texttt{conf/}; the
formulation note (\texttt{docs/FORMULATION.md}) and calibration note
(\texttt{docs/BASM\_CALIBRATION.md}) are versioned alongside the code.

\subsection{Default-$\varepsilon$ Benchmark}
Table~\ref{tab:default-grid} reports mean $F_1$, mean wall-clock
runtime, and feasibility rate per (city, $N$, solver) cell with
$\bar{R} = +\infty$ and $\bar{H} = +\infty$ (no $\varepsilon$
constraints active). \TBD{populate from data/results once the iDRAC
remote run completes; the script
\texttt{scripts/run\_grid.py} is invoked with
\texttt{--n-seeds 10 --solvers ga,alns,milp --milp-max-N 50}}.

\begin{table*}
\caption{\textbf{Default-$\varepsilon$ benchmark: mean $F_1$ (lower is better),
mean wall-clock seconds, and feasibility rate over 10 seeds.}}
\label{tab:default-grid}
\centering
\begin{tabular}{lcccccc}
\toprule
\multirow{2}{*}{City} & \multicolumn{2}{c}{$N=20$} & \multicolumn{2}{c}{$N=50$} & \multicolumn{2}{c}{$N=100$} \\
& MILP & GA & MILP & GA & ALNS & GA \\
\midrule
Bengaluru & \TBD{} & \TBD{} & \TBD{} & \TBD{} & \TBD{} & \TBD{} \\
Delhi & \TBD{} & \TBD{} & \TBD{} & \TBD{} & \TBD{} & \TBD{} \\
Gurugram & \TBD{} & \TBD{} & \TBD{} & \TBD{} & \TBD{} & \TBD{} \\
Mumbai & \TBD{} & \TBD{} & \TBD{} & \TBD{} & \TBD{} & \TBD{} \\
Pune & \TBD{} & \TBD{} & \TBD{} & \TBD{} & \TBD{} & \TBD{} \\
\bottomrule
\end{tabular}
\end{table*}

\subsection{Pareto Frontier: Risk vs.\ $F_1$}
Figure~\ref{fig:pareto} traces the $\bar{R}$-sweep at $N = 50$ for one
seed per city. As $\bar{R}$ tightens from $\infty$ down to $0.05$ the
$F_1$ cost rises monotonically and the per-route survival probability
$e^{-\bar{R}}$ correspondingly rises from \TBD{} to \TBD{}. The
operator-facing reading: a one-percentage-point increase in mean
per-route survival costs roughly \TBD{} units of $F_1$ on average across
the five cities.

\begin{figure}
\centerline{\framebox[\columnwidth]{\rule{0pt}{4cm}\textbf{[Pareto plot --- regenerate from data/results]}}}
\caption{Risk--$F_1$ Pareto frontier at $N = 50$. \TBD{regenerate
once data/results contains the iDRAC sweep output.}}
\label{fig:pareto}
\end{figure}

\subsection{Residential-Road Budget Sensitivity}
Tightening the per-route residential cap $\bar{H}_k$ from 100 to 50
edges costs \TBD{}\% in mean $F_1$ at $N = 50$ in Bengaluru and \TBD{}\%
in Mumbai, with the gap explained by Mumbai's larger residential
penetration coefficient \TBD{}.

\subsection{Solver Statistical Comparison}
For each (city, $N$) cell, we compute the Wilcoxon signed-rank
statistic \cite{wilcoxon1945} between paired GA and ALNS $F_1$ values
across the 10 seeds, applying Bonferroni correction across the
resulting
$5 \times 4 = 20$ comparisons. \TBD{report the count of cells in
which ALNS dominates GA at $\alpha = 0.05$ Bonferroni-corrected;
expected pattern based on pilot runs is that ALNS dominates at
$N \geq 100$.}

\subsection{Behavioural Compliance Sweep}
Figure~\ref{fig:behav} shows the mean per-rider survival as the
compliance prior shifts from $(\alpha, \beta) = (8, 1)$ (high
compliance) to $(3, 3)$ (low compliance). The dashed band is the IQR
across the nine operating points. \TBD{populate from data/results
once \texttt{scripts/run\_grid.py} has run with behavioural sensitivity
enabled.}

\begin{figure}
\centerline{\framebox[\columnwidth]{\rule{0pt}{4cm}\textbf{[Behavioural sensitivity --- regenerate from data/results]}}}
\caption{Behavioural-compliance sensitivity: per-rider survival vs.\
compliance prior parameters. \TBD{regenerate.}}
\label{fig:behav}
\end{figure}

\section{Discussion}
\label{sec:discussion}

\subsection{Operator-facing implications}
The $\varepsilon$-constraint formulation lets a planner read the
safety/efficiency trade-off as a tangible knob. Rather than tuning a
weighted-sum coefficient until ``it looks reasonable,'' the planner
picks a target survival probability $e^{-\bar{R}}$ and the solver
returns the cheapest fleet plan that meets it (or reports
infeasibility). The fleet-wide residential-edge budget $\bar{H}$
gives a separate dial for political/community concerns about
vulnerable-street penetration, which is operationally distinct from
per-rider safety.

\subsection{Limitations}
Three honest limitations remain. First, the BASM uses public aggregate
data and an OSM structural proxy; it does not have access to the
protected iRAD microdata, so the per-edge probability is a calibrated
prior, not a measurement. Second, the dark-store snapshot
cross-validation against OpenStreetMap returned $\sim$8\% match at
1~km radius; this reflects sparse OSM brand tagging in India rather
than a KML accuracy problem, but a manual sample audit on Google Maps
would lift this from a coverage statistic to an accuracy figure.
Third, MILP scaling to $N \geq 100$ is impractical with CBC; a
commercial solver (Gurobi/CPLEX) would close some of that gap, but
the heuristics already produce competitive bounds.

\subsection{Threats to validity}
The customer placement model uses residential-density weighting on
the OSM \texttt{landuse} layer rather than a population-density
raster. Where census or WorldPop microdata is available, plugging it
into \texttt{instance.generator} would tighten the demand model. The
behavioural compliance prior is a single field-data citation, not a
multi-study meta-analysis; the sensitivity sweep is the safeguard.

\section{Conclusion}
\label{sec:conclusion}

We presented an SA-VRPTW formulation for q-commerce last-mile
delivery in Indian cities that handles crash safety as an explicit
$\varepsilon$-constraint, calibrates per-edge crash probability
through a documented public-source chain, and is implemented under a
unified solver/evaluation contract across MILP, GA, and ALNS. The
empirical study on real Blinkit dark-store coordinates traces the
Pareto frontier between $F_1$ delivery cost and per-route crash
survival, and characterises the additional cost of network-wide
residential-road budgets. The full pipeline is publicly released
under the MIT licence.

Future work will lift three constraints. First, replacing the
behavioural compliance prior with a multi-study meta-analysis. Second,
incorporating real-time congestion feeds where available, replacing
the BPR static model. Third, extending the formulation to a stochastic
demand variant in which order arrivals across the dispatch window are
modelled by a non-stationary point process rather than a uniform
$[0, 5]$~minute draw.

\section*{Acknowledgment}
The authors thank the maintainers of \texttt{darkstoremap.in} for
publishing the dark-store KML, the OSMnx and OpenStreetMap
communities for the road-network data, and CRRI New Delhi for the
urban-flow surveys that anchor the BPR exposure assumptions.

\bibliographystyle{IEEEtran}
\bibliography{references}

\begin{IEEEbiographynophoto}{Subhajit Sidhanta}
is an Assistant Professor (Grade-I) with the Department of Industrial
and Systems Engineering, Indian Institute of Technology Kharagpur,
West Bengal, India. His research interests include distributed
systems, applied operations research, and computational
infrastructure for data-intensive urban-logistics problems.
\end{IEEEbiographynophoto}

\begin{IEEEbiographynophoto}{Pranav Bhadane}
is pursuing the Dual Degree (B.Tech.\ + M.Tech.) in the Department of
Industrial and Systems Engineering, Indian Institute of Technology
Kharagpur, West Bengal, India. His research interests include vehicle
routing, multi-objective optimisation under safety constraints, and
reproducible data pipelines for urban-logistics research.
\end{IEEEbiographynophoto}

\EOD

\end{document}
